% implications.tex -- draft Discussion content: the actionable implications,
% deduplicated against comparison-with-baselines.tex. Everything that restates
% the baseline-comparison reflection is cited there, not repeated here.
% Within each subsection, paragraphs are ordered essential, then important,
% then interesting.
%
% Integration into paper.tex:
%   1. "Implications for practitioners" is a NEW subsection. Insert it before
%      the existing "Implications for harness developers".
%   2. "Implications for AI&SE researchers" REPLACES the existing subsection of
%      the same name. The replacement keeps the normalization and harness-
%      diversity points of the existing text, and drops the two sentences that
%      rest on developer statements (interaction logs, near-zero manual work):
%      per the current rule, outreach material stays in External verification.
%   3. The provenance paragraph EXTENDS the existing "Implications for harness
%      developers" subsection (append at its end). It is a proposal, not a
%      finding.
%
% TODO before submission: align the abstract's rest-hours claim with the RQ1
% results text (the abstract says agents keep working while the developer does
% not; the results say "we assume" and the hourly data is in UTC).

\subsection{Implications for practitioners}
\label{sec:implications-practitioners}

\textbf{Set output expectations from the developer--agent system, not from
pre-2024 envelopes.}
The cohort sits outside every historical envelope, including the prolific rows,
the 99th percentile, and the plausibility ceiling
(Section~\ref{sec:baseline-comparison}).
The unit of production is now the developer--agent system, not the developer.
Teams should not set output expectations or capacity plans from pre-2024
per-developer baselines.

\textbf{Measure net output, not gross activity.}
Most file changes are changes to a file the developer changed before: 50.8\% to
88.1\% of them (\autoref{tab:rq4-rework}).
For one developer, agent state files make up 51.7\% of the churn, and only 30.7\%
of the cohort's churn is code at all (Section~\ref{sec:rq4-results}).
Gross commit and churn counts therefore make net progress look bigger than it
is.
Teams should measure net output: distinct files, net lines, and finished
milestones.
This correction runs against the contaminated baselines of
Section~\ref{sec:baseline-comparison}---one inflates the measured gap, the other
deflates it---so neither yields a multiplier.

\textbf{Read activity metrics as system output, not as the person.}
Commit frequency, work schedules, and repository switching now measure the
developer--agent system (Section~\ref{sec:baseline-comparison}).
The median commit stays small at every volume: one or two files, 27 to 88 lines
(\autoref{tab:rq4-churn-scattering}), and extreme volume and small commits
coexist.
Contribution graphs therefore track pipeline throughput, not human work, and
managers should not use commit counts as a measure of effort.
Commits also arrive during the assumed rest hours, and the hourly data is in
UTC, so commit telemetry is not a measure of presence.

\textbf{Plan on sustained rates, not on peaks.}
Burst-collapse is the modal trajectory, and yearly baselines smooth such bursts
away (Section~\ref{sec:baseline-comparison}).
Teams should set quotas and capacity plans on sustained rates, not on peak
months.

\textbf{Track total debt, not only per-line rates.}
``No quality degradation'' is a per-line result (Section~\ref{sec:rq5-results}):
the code grows faster than the per-line rates fall, so total remediation effort
grows; duplication worsens in 60\% of projects; and the share of code in files
over 1{,}000 lines doubles, from 22\% to 46\%.
Teams should track total debt and the file-size distribution, not only per-line
rates.
This extends the existing code-quality implications
(Section~\ref{sec:implications-code-quality}).

\textbf{Hold agents to small-commit discipline.}
The typical commit touches one or two files at every volume, and between 76.4\%
and 89.0\% of commits touch at most five files (\autoref{tab:rq4-breadth}).
Teams can set rules for small, focused commits and still keep high volume.

\textbf{Govern agent-written state.}
Agents write their own coordination state into the repository: in-repository
issue trackers, agent memory, and agent instructions
(Section~\ref{sec:rq4-results}).
For one developer these files are 15.4\% of file changes and 51.7\% of churn.
They make activity counts too high, and they fill review queues.
Teams should decide what agent state lives in the repository, and keep the rest
out.

\textbf{Treat repository breadth as a choice.}
The cohort keeps 7.2 active repositories per week against a pre-2024 baseline of
3.7, and the median developer switches repositories 4 times per day (RQ3).
Yet most effort sits in a few repositories: the median weekly Gini coefficient
is 0.449.
Teams should treat breadth as a choice, and expect per-project velocity to look
modest even when aggregate output is extreme.

\subsection{Implications for AI\&SE researchers}
\label{sec:implications-researchers}

\textbf{Treat agent-detection signals as a lower bound.}
A commit without an agent signal may still be agent work: three of the nine
developers show near-zero signals, two of them at extreme commit volumes
(Section~\ref{sec:rq2-results}).
Non-attributed commits may include work by agents that leave no detectable
signals, as well as genuine manual work.
Commit-trace studies therefore give lower bounds for agent use.
Researchers should move away from relying on deliberate harness disclosure in
favor of more robust signals, and should combine commit signals with
independent signals.
Adoption measurements should also normalize per developer: the unnormalized
harness distribution is heavily skewed by a few highly prolific developers
(Section~\ref{sec:rq2-results}).
Detectable harness diversity is limited as well, because the harness
distribution of the actual implementation is abstracted away when code is
committed.

\textbf{Filter churn by file category before comparing repositories.}
Code is 30.7\% of the cohort's churn; the residual category is 28.2\% and holds
million-line machine artifacts (Section~\ref{sec:rq4-results}).
Filters written before 2025 do not know agent state files.
Researchers should filter churn by file category before comparing repositories
across eras.

\textbf{Build the missing baseline: developers with agents at ordinary
volumes.}
The baseline comparison yields a bound, not a multiplier, because the
populations and the eras differ (Section~\ref{sec:baseline-comparison}).
The missing baseline is developers with agents at ordinary volumes.
Only that study would turn the bound into a causal estimate.
Researchers should build that baseline; it is also the natural next step for
this work (Section~\ref{sec:conclusion}).

\textbf{Do not classify work as new or rework from first touch.}
Same-day revisits are 23.9\% to 42.5\% of changes, and a file receives 2.03 to
8.42 changes over the window (\autoref{tab:rq4-rework}).
First-touch classifications therefore misclassify agent iteration.

\textbf{Use the developer as the unit.}
Trajectories diverge between sustained growth and burst-collapse, and
cross-sectional means hide both (Section~\ref{sec:results}).
Researchers should group developers by trajectory, not by cross-sectional means.

\textbf{Separate accumulation from aging, and look beyond file-local metrics.}
Our design separates code accumulation from aging after construction, and all
five quality metrics are file-local (Section~\ref{sec:rq5-results}); none of
them describes the codebase as a whole.
Researchers should add measures such as change locality and coupling.

% Append to the existing "Implications for harness developers" subsection:

\textbf{Provenance must become a protocol property, not a trailer. (Proposal)}
A commit trailer is a text string, and nothing stops a tool from omitting it or
from appending it to work it did not do.
Where provenance is a compliance or credit requirement, no trailer-based
pipeline satisfies it.
Harness builders should add a provenance mechanism that tools cannot omit or
forge.
This is a proposal for tool builders, not a finding of this study.
